% stable-fips203-mlkem-pke-decrypt-k4 — FIPS 203 Alg.15 ML-KEM-1024 PKE Decrypt (k=4)
% 编译：bash ../../../scripts/xelatex-clean.sh stable-fips203-mlkem-pke-decrypt-k4-实现方案-customspec.tex
\documentclass[UTF8,a4paper,11pt]{ctexart}
\usepackage{amsmath,amssymb}
\usepackage{booktabs}
\usepackage{geometry}
\usepackage{hyperref}
\usepackage{longtable}
\usepackage{array}
\usepackage{seqsplit}
\geometry{margin=2.2cm}
\newcolumntype{L}[1]{>{\raggedright\arraybackslash}p{#1}}
\newcommand{\apiname}[1]{\texttt{\seqsplit{#1}}}
\newcommand{\dirname}[1]{\texttt{\seqsplit{#1}}}
\hypersetup{colorlinks=true,linkcolor=blue,urlcolor=blue}

\title{stable-fips203-mlkem-pke-decrypt-k4\\FIPS 203 Alg.15 ML-KEM-1024 PKE Decrypt 实现方案（k=4）}
\author{ascendc 工程（定型交付）}
\date{2026-07-09}

\begin{document}
\maketitle

\begin{abstract}
本文档描述 \textbf{定型交付算子} \dirname{examples/stable/ml-kem/ml-kem-1024/stable-fips203-mlkem-pke-decrypt-k4/}：
在 \textbf{Atlas A2 / Ascend910B4} 上以 AscendC 实现 FIPS 203 \textbf{Algorithm 15}（ML-KEM-1024，$K=4$）\textbf{K-PKE.Decrypt} 全链。

\textbf{生产 I/O（强制对齐 Alg.15）}：
\begin{itemize}
  \item \textbf{输入}：\texttt{dk\_PKE}（1536\,B）+ 密文 $c=c_1\Vert c_2$（1568\,B）；另加与 seed 无关的\textbf{静态} NTT/INTT LUT（非本算法中间态）。
  \item \textbf{输出}：\textbf{仅}消息 $m$（32\,B）。
  \item \textbf{禁止}将 $u'$、$v'$、$\hat{s}$、$\hat{u}$、$\hat{w}$、$w$ 等中间量以 \texttt{.bin} 落盘到 \texttt{input/}/\texttt{output/}，或经 Host D2H/H2D 中转冒充全链。
\end{itemize}

\textbf{自包含原则}：AscendC 与 Host golden 源码 \textbf{全部 vendored 于本目录}。
\textbf{唯一允许的外部代码依赖}：仓库 \dirname{library/shared/}。
\textbf{禁止}：运行时 \texttt{\#include}/import \dirname{ascendc-tests/}、其它 \dirname{examples/}、\dirname{frozen/}、liboqs（KAT 脚本除外）。

\textbf{设备 Launch}（与 PASS 探针 \dirname{pass-fix-f203-alg15-pke-decrypt-device-k4} 同构）：
SIM/CPU \textbf{1} 次 \texttt{f203\_decrypt\_device\_fused}（\texttt{KERNEL\_TYPE\_MIX\_AIC\_1\_2}，$\texttt{blockDim=1}$）。
中间 GM \textbf{仅设备内}；生产路径\textbf{仅} D2H $m$。
SIM 910B4 参考 Total tick $\approx$ \textbf{283k}。

技术基线：\dirname{docs/notes/F203-Alg15-Decrypt-2launch编排技术总结.md}（历史 2-launch 约束；本规格采用已验证的 GATE 4/8 单 kernel）、
\dirname{docs/notes/F203-Compress-Decompress-向量实现指南.md}、
\dirname{docs/notes/F203-ByteEncode-ByteDecode-d-向量与标量选型.md}、
\dirname{docs/engineering/用例自包含与设备全链约束.md}。
写法对齐 \dirname{exp-fips203-mlkem-pke-encrypt-k4-实现方案-customspec.tex}。
\end{abstract}

\tableofcontents
\newpage

\section{工程约束（自包含）}

\subsection{允许的外部依赖}

\begin{longtable}{@{}L{4.5cm}L{8.5cm}@{}}
\toprule
路径 & 用途 \\
\midrule
\dirname{library/shared/} & 公共头（若 Host/设备派生需要） \\
CANN / tikicpulib & 工具链与仿真（非业务源码） \\
\bottomrule
\end{longtable}

\textbf{说明}：Decrypt \textbf{无} SampleNTT / CBD / PRF，一般\textbf{不}依赖 SHAKE 内嵌；若 vendor 路径误带入，须切断。

\subsection{禁止的外部依赖与不安全 I/O}

\begin{itemize}
  \item \textbf{ascendc-tests} 任意探针目录的编译期 include / Python import（允许\textbf{一次性 vendor 抄入}本目录后切断外链）。
  \item \textbf{frozen/}、G4 correctness 标量尾路线、liboqs 作为生产运行时依赖。
  \item 将 $u'/v'/\hat{s}/\hat{u}/\hat{w}/w$ 等写成 \texttt{input/*.bin} 或 \texttt{output/*.bin} 作为全链契约。
  \item Host 对拍以外的密码学计算冒充设备输出（如 Host 算 $m$ 写回 \texttt{output/m.bin}）。
  \item 段间 \textbf{D2H 再 H2D} 中转中间态。
  \item 同 kernel 内\textbf{标量写 GM} 再 MTE 读（Encrypt R2 / Decrypt pad/ŝ 教训）：须 UB + \apiname{DataCopy}/\apiname{Duplicate}。
\end{itemize}

\subsection{唯一交付路径}

本用例 \textbf{仅} 下列生产路径；\texttt{run.sh} 默认即全量，\textbf{禁止}要求用户手动 \texttt{export SIM\_DIRECT=1} 或一长串与 default 相同的宏：

\begin{longtable}{@{}L{4.2cm}L{8.8cm}@{}}
\toprule
组件 & 锁定配置 \\
\midrule
Prep unpack & ByteDecode$_{11/5/12}$ \textbf{标量} + Decompress$_{11/5}$ \textbf{向量}；落盘经 UB+\apiname{DataCopy} \\
NTT($u'$) / INTT($\hat{w}$) & poly-batch；平面 mat\_c；S1--S3 \textbf{禁止 Gather} \\
su\_dot & $\hat{w}\leftarrow\sum_j\mathrm{MultiplyNTTs}(\hat{s}[j],\hat{u}[j])$；AIV0 独占 + softSyncGm \\
尾段 & $v'-w$ 向量 wrap\_mod $\to$ Compress$_1$ \textbf{向量 Barrett 原地} $\to$ ByteEncode$_1$ \textbf{标量}按字节拼 \\
Host & \texttt{main}：\textbf{1} launch；\texttt{dk}+$c$+LUT $\to$ \textbf{仅} $m$ \\
\bottomrule
\end{longtable}

\subsection{目录布局（计划；写码阶段落地）}

\begin{verbatim}
stable-fips203-mlkem-pke-decrypt-k4/
  stable-fips203-mlkem-pke-decrypt-k4-实现方案-customspec.tex/.pdf  # 本规格
  main.cpp / main_decrypt_*.cpp    # 生产 Host：1 launch
  f203_decrypt_layout.h            # GM 布局
  compute/g4_full/                 # f203_decrypt_device_fused
  compute/ntt_u/ intt_w/ su_dot/   # vendored NTT/INTT/内积
  compute/compress1_byteencode1/   # 尾融合
  unpack/ 或 prep/                 # Decode + Decompress
  CMakeLists.txt / cmake/
  scripts/gen_data.py              # dk/c/lut + golden_m
  scripts/host_golden/
  run.sh                           # 默认 SIM_DIRECT=1
  SELF_CONTAINED.md                # 写码时补
  STATUS.md                        # 写码验收后补
\end{verbatim}

\textbf{参考实现过程（只读，写码时 vendor，禁止运行时外链）}：
\dirname{ascendc-tests/ml-kem/ml-kem-1024/pass-fix-f203-alg15-pke-decrypt-device-k4/}。

\section{数学语义（Alg.15，$K=4$，$N=256$，$q=3329$）}
\label{sec:math}

\subsection{参数与符号}

\begin{itemize}
  \item 参数集：\textbf{ML-KEM-1024}（FIPS 203；$k=4$；$d_u=11$；$d_v=5$）。
  \item $c_1$：$k\cdot 352=1408$\,B；$c_2$：$160$\,B；$|c|=1568$\,B。
  \item $\texttt{dk\_PKE}$：$384k=1536$\,B（$\mathrm{ByteEncode}_{12}(\hat{s})$）。
  \item \textbf{无采样}：无 SampleNTT / CBD / PRF / coins；\textbf{无矩阵} $\hat{A}$。
\end{itemize}

\subsection{输入输出（生产契约）}

\begin{longtable}{@{}L{3.5cm}L{2.2cm}L{7.5cm}@{}}
\toprule
路径 & 尺寸 & 说明 \\
\midrule
\texttt{input/dk\_pke.bin} & 1536\,B & 私钥多项式编码 \\
\texttt{input/c.bin} & 1568\,B & $c=c_1\Vert c_2$ \\
\texttt{input/lut\_*.bin} & $4\times 64$\,KB & \textbf{静态} NTT/INTT even/odd planar-stacked；与 seed 无关 \\
\texttt{output/m.bin} & \textbf{32\,B} & \textbf{唯一产物} \\
\bottomrule
\end{longtable}

调试 Gate（须显式 env，\textbf{非默认}）可 D2H 中间量；生产 \texttt{DECRYPT\_GATE=0}。

\subsection{流水线（FIPS 行 1--8）}

\begin{enumerate}
  \item 行 1--2：$c_1\leftarrow c[0:1408]$；$c_2\leftarrow c[1408:1568]$。
  \item 行 3：$u'\leftarrow\mathrm{Decompress}_{11}(\mathrm{ByteDecode}_{11}(c_1))$（$k$ poly）。
  \item 行 4：$v'\leftarrow\mathrm{Decompress}_{5}(\mathrm{ByteDecode}_{5}(c_2))$（1 poly）。
  \item 行 5：$\hat{s}\leftarrow\mathrm{ByteDecode}_{12}(\texttt{dk\_PKE})$（$k$ poly）。
  \item 行 6：$w\leftarrow v'-\mathrm{NTT}^{-1}\bigl(\hat{s}^\top\circ\mathrm{NTT}(u')\bigr)$。
  \begin{itemize}
    \item $\hat{u}\leftarrow\mathrm{NTT}(u')$；
    \item $\hat{w}\leftarrow\sum_{j=0}^{k-1}\mathrm{MultiplyNTTs}(\hat{s}[j],\hat{u}[j])$；
    \item $w\leftarrow v'-\mathrm{INTT}(\hat{w}_{\mathrm{padded}})$（$\hat{w}$ 前缀 pad 至 $k$ poly，余零）。
  \end{itemize}
  \item 行 7：$m\leftarrow\mathrm{ByteEncode}_{1}(\mathrm{Compress}_{1}(w))$（32\,B）。
  \item 行 8：return $m$。
\end{enumerate}

\textbf{行 7 方向}：是 $\mathrm{ByteEncode}_1(\mathrm{Compress}_1(w))$，\textbf{不是} $\mathrm{ByteDecode}_1$，也不是 Compress「逆」。
\textbf{禁止}把 Encrypt 行 20 的 $\mu\leftarrow\mathrm{Decompress}_1(\mathrm{ByteDecode}_1(m))$ 当 Decrypt 尾段抄。

\textbf{本阶段不做}：KEM Decaps（Alg.21）、KeyGen、Encrypt；Decode$_{11/5/12}$ 真·向量 unpack（低 ROI）；su\_dot$\to$INTT UB 驻留（两档实验已否决，见 §\ref{sec:closed}）。

\subsection{Compress$_1$ / ByteEncode$_1$（定点）}

\begin{itemize}
  \item Compress$_1$（liboqs Barrett，与 golden 对拍）：
  \[
    \mathrm{bit}=\bigl((u\cdot 1290168)+(1\ll 30)\bigr)\gg 31,\quad u\in[0,q).
  \]
  设备：\apiname{Muls}/\apiname{Adds}/\apiname{ShiftRight}；int32 环绕 $\equiv$ u32 Barrett。
  \item ByteEncode$_1$：LSB-first 标量 pack，$256\times 1\,\mathrm{bit}\to 32$\,B；\textbf{不做}真·向量 bit 流。
\end{itemize}

\subsection{Golden（I/O 等价）}

\begin{itemize}
  \item 锁死种子 \texttt{SEED\_D=20260619}（与 Encrypt / round-trip 同源）。
  \item Host：\texttt{scripts/host\_golden/golden\_m.py}（自包含）$\to$\texttt{output/golden\_m.bin}。
  \item 验收：\texttt{output/m.bin} vs \texttt{golden\_m.bin} \textbf{max\_abs\_diff=0}（CPU $\land$ SIM）。
  \item \textbf{禁止}以「与探针源码同构」为验收；仅 I/O 等价。
\end{itemize}

\section{编排与 Launch}
\label{sec:launch}

\subsection{SIM / CPU 生产路径（单 kernel）}

\begin{verbatim}
Session 1 aclrtStream
  Launch 1  f203_decrypt_device_fused   blockDim=1  MIX_AIC_1_2
            in:  dk_pke, c, ntt_ws(+LUT), intt_ws(+LUT), softSyncGm(=0)
            GM:  u,v,s_hat,u_hat,w_hat,w_padded,w_time（设备内，不 D2H）
            out: m[32]  # 唯一 D2H
\end{verbatim}

合计 \textbf{1 launch}；参考 tick $\approx$ \textbf{283k}。

\subsection{核内 FSM（锁定）}

\begin{longtable}{@{}L{3.5cm}L{9.5cm}@{}}
\toprule
段 & 同步 \\
\midrule
prep（unpack+decode $\hat{s}$） & 仅 AIV0；AIV1 经 softSyncGm slot0 等待；双 AIV \texttt{Set(4)}；AIC \texttt{Wait(4)}/\texttt{Set(8)} \\
NTT($u'$) & CrossCore flag \textbf{1/2/3}（SPLIT/MMAD/PACK） \\
su\_dot + pad\_w\_hat & 仅 AIV0；AIV1 softSyncGm slot1；再 GATE 4/8 \\
INTT($\hat{w}$) & flag \textbf{1/3}（\textbf{禁}复用 2）；GATE 后尾段 \\
尾 $v'-w$ / Compress$_1$ / Encode$_1$ & AIV 向量 + 标量 Encode；写 $m$ \\
\bottomrule
\end{longtable}

\textbf{历史对照}：G4 correctness 为 prep $\Vert$ NTT+INTT \textbf{2 launch}（见 2-launch 笔记）。
单 kernel 可行条件：段间用 \textbf{GATE flag 4/8} 分隔 NTT 与 INTT，避免同 flag 1--3 争用；prep/su\_dot 用 softSyncGm 汇合双 AIV。

\subsection{su\_dot $\to$ INTT（生产）}

\begin{itemize}
  \item $\hat{w}$ 写 \texttt{wHatGm}；\texttt{pad\_w\_hat\_for\_intt}：UB \apiname{Duplicate} 零填充 + \apiname{DataCopy} $\to$ \texttt{wPaddedGm}；INTT Stage1 \texttt{Process()} 读 GM。
  \item \textbf{已否决}：UB 驻留 / \texttt{ProcessFromLocal}（半吊子与同 TPipe 两档，SIM $\sim$287k，高于基线 $\sim$283k）。
\end{itemize}

\section{AI Core 规模}
\label{sec:aicore}

\begin{longtable}{@{}L{3.2cm}L{10cm}@{}}
\toprule
项 & 锁定 \\
\midrule
核类型 & \texttt{KERNEL\_TYPE\_MIX\_AIC\_1\_2} \\
\texttt{blockDim} & \textbf{1}（\texttt{aicore=1}） \\
角色 & \textbf{1 AIC + 2 AIV}：AIC 做 NTT/INTT Stage2 MMAD；AIV 做 prep/split/pack/su\_dot/尾 \\
数据划分 & $k{=}4$ poly-batch；每 AIV 握完整 poly 的 hi+lo（\textbf{禁} limbsplit） \\
选型理由 & 与 PASS 探针同构；Decrypt 无 SampleNTT 双核 prep 需求 \\
\bottomrule
\end{longtable}

\textbf{禁止}实现擅自改 \texttt{blockDim} 却不回写本规格。
串行占用：1 颗 AI Core（Cube+2 Vector）。CPU \texttt{[SUCCESS][AIC\_*]} 为 tikicpu artifact；以 SIM Total tick 为准。

\section{数据与组件同步规划}
\label{sec:sync}

\begin{itemize}
  \item Host：H2D \texttt{dk}/$c$/LUT/softSync；\textbf{一次} launch；D2H \textbf{仅} $m$。
  \item 核内：CrossCore GATE 4/8 + softSyncGm（AIV0/1）；\apiname{PipeBarrier} 保证 UB/GM 可见性。
  \item 调试：\texttt{F203\_DECRYPT\_TRACE=1} 时 Host 分配独立 \texttt{traceGm} 并轮询段标记（非默认）。
  \item GM 可见性：\texttt{decode\_s\_hat} / \texttt{pad\_w\_hat} \textbf{禁止}标量写 GM 再 MTE 读。
\end{itemize}

\subsection{实现映射（写码时）}

\begin{verbatim}
main.cpp
  -> Launch1: f203_decrypt_device_fused (prep|NTT|dot|INTT|tail)
  -> D2H: m only

f203_decrypt_layout.h
  -> static_assert 各 GM 缓冲尺寸
\end{verbatim}

\section{全量 API 清单（查阅 CANN 9.0.0）}
\label{sec:api-table}

说明：分类按 PDF 第 2 章；在线页为社区 \texttt{atlasascendc\_api\_07\_*}。
细粒度算子已在上游 PASS 探针验证；写码以 vendor 后本目录源码为准，并遵守 Rule「写 AscendC 前查索引」。

\begin{longtable}{@{}L{2.8cm}L{1.6cm}L{2.2cm}L{0.7cm}L{3.2cm}L{3.2cm}@{}}
\toprule
API & 分类 & 在线文档 ID & A2 & 输入/输出 dtype & 备注 \\
\midrule
\endhead
\apiname{GlobalTensor} & 2.2 & 07\_0007 & $\checkmark$ & GM 描述 & \texttt{SetGlobalBuffer} \\
\apiname{LocalTensor} & 2.2 & 07\_0006 & $\checkmark$ & UB 描述 & Alloc/Get \\
\apiname{GetBlockIdx}/\apiname{GetSubBlockIdx} & 2.3 系统 & 07\_0185 邻域 & $\checkmark$ & 核/子核索引 & MIX 角色；AIV0 独占段 \\
\apiname{KERNEL\_TASK\_TYPE\_DEFAULT} & Utils & 编程指南 & $\checkmark$ & MIX\_AIC\_1\_2 & 单核类型 \\
\apiname{TPipe}/\apiname{TQue}/\apiname{TBuf} & 2.3 资源 & 07\_0108/0137 & $\checkmark$ & UB 管道 & 向量段 \\
\apiname{DataCopy} & 2.3.1 & 07\_0101 & $\checkmark$ & int32 GM$\leftrightarrow$UB & pad/ŝ/尾；对齐约束 \\
\apiname{Duplicate} & 2.3.3 & 查阅索引 & $\checkmark$ & int32 & pad 零填充 \\
\apiname{PipeBarrier} & 2.3 同步 & 查阅索引 & $\checkmark$ & — & UB/GM 可见性 \\
\apiname{CrossCoreSetFlag}/\apiname{CrossCoreWaitFlag} & 跨核 & API 列表 & $\checkmark$ & MIX AIC$\leftrightarrow$AIV & NTT 1/2/3；INTT 1/3；GATE 4/8 \\
\apiname{Add}/\apiname{Sub}/\apiname{Muls}/\apiname{Mul} & 2.3.3 & 07\_0035--0055 & $\checkmark$ & int32 & su\_dot / Compress$_1$ / wrap\_mod \\
\apiname{Max} & 2.3.3 & 查阅索引 & $\checkmark$ & int32 & $v'-w$ wrap\_mod \\
\apiname{ShiftRight}/\apiname{ShiftLeft} & 2.3.3 & 07\_0059 & $\checkmark$ & int32 & Compress$_1$/Decompress \\
\apiname{Cast} & 2.3.3 & 07\_0073 & $\checkmark$ & 见 Cast 链 & 按需；A2 无直连 int32$\to$int8 \\
\apiname{Gather} & 2.3.3 & 07\_0071 邻域 & $\checkmark$ & — & \textbf{禁止} Tag5T S1--S3；尾/ basemul 外另议 \\
Cube MMAD / \apiname{AicMmad} 路径 & 2.3.2 矩阵 & 教程+工程封装 & $\checkmark$ & int8$\times$int8$\to$int32 & Stage2；B 为 ws 内 LUT \\
\bottomrule
\end{longtable}

\subsection{评估过但未采用 / 已关闭}
\label{sec:closed}

\begin{longtable}{@{}L{4.5cm}L{8.5cm}@{}}
\toprule
项 & 原因 \\
\midrule
su\_dot$\to$INTT UB 驻留（半吊子 / 同 TPipe） & SIM $\sim$287k $>$ 基线 $\sim$283k；CPU \texttt{ProcessFromLocal} 对拍失败 \\
ByteEncode$_1$ 真·向量 bit 流 & ROI 差；选型标量 pack \\
ByteDecode$_{11/5/12}$ 真·向量 & 低 ROI；prep 占比小 \\
NTT S1--S3 内 \apiname{Gather} & MLKEM Tag5T 禁令 \\
G4 2-launch 作为本 exp 生产路径 & 已被单 kernel PASS 取代；correctness 探针仅作语义对照 \\
中间态 bin staging & 违反 Alg.15 I/O \\
\bottomrule
\end{longtable}

\section{数学步骤 $\to$ API 覆盖率}
\label{sec:coverage}

\begin{longtable}{@{}L{4.2cm}L{4.0cm}L{1.5cm}L{3.5cm}@{}}
\toprule
数学步骤 & 主路径 & 覆盖 & 缺口/备注 \\
\midrule
ByteDecode$_{11/5/12}$ & \textbf{标量} unpack + DataCopy & 部分 & 有意标量；非隐瞒 \\
Decompress$_{11/5}$ & 向量 Muls/Adds/ShiftRight & 高 & 对齐 decompress 探针 \\
NTT($u'$) S1--S3 & 向量 split + AIC MMAD + pack & 高 & 无 Gather；poly-batch \\
su\_dot MultiplyNTTs & 向量/UB（AIV0） & 高 & softSync 汇合 \\
pad $\hat{w}$ & Duplicate+DataCopy & 100\% & 禁标量写 GM \\
INTT($\hat{w}$) & MIX flag 1/3 & 高 & GATE 4/8 分隔 \\
$v'-w$ wrap\_mod & Sub/Max/ShiftRight & 高 & 向量 \\
Compress$_1$ & 向量 Barrett & 高 & 原地写回 \\
ByteEncode$_1$ & \textbf{标量} pack & 部分 & 有意标量 \\
$m$ 写出 & DataCopy GM & 100\% & 唯一 D2H \\
\bottomrule
\end{longtable}

主路径：\textbf{非}宣称全链 100\% 向量（Decode$_{*}$ 与 Encode$_1$ 标量已定型）；NTT/INTT/su\_dot/Compress$_1$/$v'-w$ 与 PASS 探针一致。

\section{验证}
\label{sec:verify}

\subsection{生产验收（默认；写码完成后）}

\begin{verbatim}
cd examples/stable/ml-kem/ml-kem-1024/stable-fips203-mlkem-pke-decrypt-k4
bash run.sh -r cpu -v Ascend910B4
bash run.sh -r sim -v Ascend910B4
# 期望：两模式 m max=0；output/ 仅 m.bin；SIM 根目录 0 stray dump
# 参考 tick ~ 283k
\end{verbatim}

\textbf{禁止}标准文档要求手动 \texttt{SIM\_DIRECT=1} 或与 default 重复的 \texttt{HAT\_*} export。

\subsection{自包含 golden}

\texttt{scripts/gen\_data.py}：\texttt{SEED\_D=20260619} 生成 \texttt{dk}/$c$/LUT/\texttt{golden\_m}（可与 KeyGen/Encrypt fixtures 同源派生）。

\subsection{可选闭环 / liboqs}

\begin{itemize}
  \item 仓库 \texttt{scripts/roundtrip\_pke\_encrypt\_decrypt.sh}：stable KeyGen+Encrypt $\to$ 本 Decrypt（写码后可将 \texttt{DECRYPT\_DIR} 指本目录）。
  \item \texttt{scripts/liboqs\_pke\_vs\_ascendc.sh}：Decrypt 段 vs liboqs（可选）。
\end{itemize}

\section{性能（SIM 910B4 参考）}

上游 PASS 探针（2026-07-09）：Total tick $\approx$ \textbf{283248--283252}（1 launch）。
本 exp 写码验收目标：与上述同量级；\textbf{远超}视为回归，不得靠拉长防挂死预算掩盖。

\section{写码衔接（本规格闭环后）}

\begin{enumerate}
  \item 用户确认本 customspec 路径后，执行 \textbf{【预研】}：按本规格 vendor
        \dirname{pass-fix-f203-alg15-pke-decrypt-device-k4} 到本目录并自包含化。
  \item 活跃 customspec：\dirname{examples/stable/ml-kem/ml-kem-1024/stable-fips203-mlkem-pke-decrypt-k4/stable-fips203-mlkem-pke-decrypt-k4-实现方案-customspec.tex}（及同名 PDF）。
  \item 晋级 \dirname{examples/stable/ml-kem/ml-kem-1024/stable-fips203-mlkem-pke-decrypt-k4} 须另走 \texttt{\#交付\#}（T15a），且须 baseline-registry 齐备。
\end{enumerate}

\end{document}
